\clearpage\fancyhead[L]{\small KOMET.}\hypertarget{native-348-lower}{}\label{native-348-lower}
\begin{center}\textbf{KOMET.}\par\vspace{4pt}\hrule\vspace{4pt}\fontsize{10}{12}\selectfont Monatliche Correspondenz. Band XXIV, S. 595\,..\,598, Dezember 1811.\par\vspace{4pt}\hrule\vspace{4pt}\end{center}

\par\indent In Göttingen fanden die Herren \textsc{Gauß} und \textsc{Harding} den Kometen gleich am Abend des 9. Dezember, an dem mein Brief dort eingegangen war. Am Kreismikrometer machte \textsc{Gauß} die folgenden Beobachtungen.

\begin{center}
\fontsize{10}{12}\selectfont\setlength{\tabcolsep}{4pt}\renewcommand{\arraystretch}{1.15}\begin{tabular}{l|c|c|c}
 & Mittlere Zeit & Rektasz. des Kometen & Südl. Dekl. \\
 & in Göttingen & & \\ \hline
Dez. 9 & $10^{\mathrm u}\;\;6^{\mathrm m}\;52^{\mathrm s}$ & $63^\circ\;49'\;41''4$ & $10^\circ\;21'\;55''5$ \\
11 & $10\;\;34\;\;\;\;1$ & $63\;\;33\;\;18.0$ & $8\;\;39\;\;46.4$ \\
12 & $\;\;8\;\;\;\;5\;\;52$ & $63\;\;26\;\;25.3$ & $7\;\;54\;\;25.9$ \\
\end{tabular}
\end{center}

\SourceLeaf{349}
\par\indent „Hier, schrieb uns \textsc{Gauß} aus Göttingen am 15. Dez. 1811, sind auch meine vorläufigen Elemente. Sie sind freilich nur als Annäherung anzusehen, zeigen aber doch, dass der Komet ein \emph{neuer} ist und an Helligkeit nicht mehr zunehmen wird.

\Needspace{16\baselineskip}
\begin{center}\fontsize{10}{12}\selectfont\setlength{\tabcolsep}{4pt}\renewcommand{\arraystretch}{1.15}\begin{tabular}{l}
Zeit des Periheldurchgangs 1811 Nov. 12,6225 Gött. Mer.\\
\end{tabular}\end{center}
\begin{center}\fontsize{10}{12}\selectfont\setlength{\tabcolsep}{3pt}\renewcommand{\arraystretch}{1.15}\begin{tabular}{llr}
Länge des Perihels & .\;.\;.\;.\;.\;.\;. & $48^\circ\;30'\;20''$\\
Log. des Abstands im Perihel & .\;.\;.\;.\;. & $0.20160$\\
Aufsteigender Knoten & .\;.\;.\;.\;.\;. & $92^\circ\;46'\;59''$\\
Neigung der Bahn & .\;.\;.\;.\;.\;.\;. & $31\;37\;55$\\
\end{tabular}\end{center}
\noindent Bewegung rechtläufig


\par\indent „Der Komet kann noch \emph{sehr} lange sichtbar sein, wenn seine Helligkeit es zulässt. Für den 31. Januar 1812, $10^{\mathrm u}$, finde ich Rektasz. $67\frac{1}{2}^\circ$, Dekl. $25\frac{1}{2}^\circ$ nördlich, Helligkeit 0.178.

\begin{center}
\begin{tabular}{l}
Helligkeit 1811 November 18 $\;=\;0.713$ \\
\qquad --- \qquad --- Dezember 12 $\;=\;0.698$ \\
\end{tabular}
\end{center}


\Needspace{20\baselineskip}\par\indent Herr \textsc{Gerling} verglich alle seither bekannt gewordenen Beobachtungen mit diesen Elementen und erhielt die folgenden Ergebnisse.

\begin{center}
\fontsize{10}{12}\selectfont\setlength{\tabcolsep}{6pt}\renewcommand{\arraystretch}{1.15}\begin{tabular}{l|c|c|l}
 & \multicolumn{2}{c|}{Abweichung} & \\
 & in Rektasz. & in Dekl. & Beoobachter \\ 
\hline
Nov. 18 & $+32''$ & $-101''$ & v. Zach \\
19 & $+80$ & $+46$ & v. Zach \\
20 & $+23$ & $+4$ & v. Zach \\
21 & $+3$ & $+15$ & v. Zach \\
Dez. 8 & $+19$ & .\;.\;.\;. & v. Lindenau \\
9 & $+11$ & $+12$ & v. Lindenau \\
9 & $+8$ & $-55$ & Olbers \\
9 & $+1$ & $+1$ & Gauss \\
11 & $-10$ & $-16$ & Gauss \\
12 & $+1$ & $+7$ & Gauss \\
\end{tabular}
\end{center}


\par\indent Um die künftigen Beobachtungen des Kometen zu erleichtern, berechnete Herr \textsc{Nicolai} aus den obigen Elementen die folgende Ephemeride.

\begin{center}
\fontsize{10}{12}\selectfont\setlength{\tabcolsep}{4pt}\renewcommand{\arraystretch}{1.15}\begin{tabular}{l|c|c|c|c}
Mitternacht & Rektasz. & Deklination & \shortstack{Log. Abst.\\$a\;\EarthSymbol$} & \shortstack{Helligkeit} \\ \hline
1811 Dezember 21 & $62^\circ\;36'$ & $0^\circ\;7$ S. & 9.8879 & 0.596 \\
29 & $62\;\;25$ & $6\;\;20$ N. & 9.9185 & 0.496 \\
1812 Januar 6 & $62\;\;49$ & $12\;\;11$ --- & 9.9562 & 0.397 \\
14 & $63\;\;47$ & $17\;\;15$ --- & 9.9984 & 0.310 \\
22 & $65\;\;17$ & $21\;\;33$ --- & 0.0427 & 0.239 \\
30 & $67\;\;17$ & $25\;\;\;\,8$ --- & 0.0872 & 0.183 \\
Februar 7 & $69\;\;41$ & $28\;\;\;\,7$ --- & 0.1310 & 0.141 \\
\end{tabular}
\end{center}


\par\indent Danach nimmt die Helligkeit des Kometen schon wieder sehr merklich ab. Da das Mondlicht seine Sichtbarkeit nun eine Zeit lang beeinträchtigen wird, bleibt es wohl noch zweifelhaft, ob sich eine weitere Reihe guter Beobachtungen dieses Kometen gewinnen lässt. Doch die schöne Übereinstimmung mit allen Beobachtungen zeigt, dass die obigen Elemente seinen Lauf schon jetzt mit sehr guter Annäherung bestimmen, sodass unsere Nachkommen ihn bei einem späteren Wiedererscheinen ohne Mühe wiedererkennen können.
